% !TeX root = proyecto.tex
%=========================================================
\chapter{Modelo de la interacción}	
\label{cap:modInteraccion}

\cdtInstrucciones{Introduzca el capítulo indicando su contenido y organización.}
Este capítulo describe el modelo de interfaces de usuario del sistema SICORRE, 
detallando la interacción entre el usuario y el sisteman. Se presentan el mapa 
de navegación general del sistema, los elementos comunes a todas las pantallas 
y la especificación detallada de cada interfaz, incluyendo sus objetivos, 
entradas, y salidas .


%---------------------------------------------------------
\section{Modelo de navegación}

La navegación del sistema SICORRE está basada en una arquitectura de tipo 
jerárquica con acceso centralizado desde un panel principal diferenciado por 
rol. Todos los flujos de navegación parten de la pantalla de inicio de sesión 
(\IUref{IU01}{Pantalla de Inicio de Sesión}), desde la cual el sistema 
redirige automáticamente al usuario hacia su panel correspondiente según 
el perfil autenticado.

El sistema cuenta con dos flujos de navegación principales:

\begin{description}
	\item[Flujo de Administrador:] Accesible para la directora y la 
	coordinadora. Desde la \IUref{IU02}{Pantalla del Panel de Administrador} 
	se puede acceder a la gestión de usuarios, equipos multidisciplinarios, 
	catálogo de actores en materia de derechos y consulta de expedientes.
	
	\item[Flujo de Especialista:] Accesible para el trabajador social, 
	abogado, psicólogo y médico. Desde la 
	\IUref{IU03}{Pantalla del Panel de Especialista} se puede acceder al 
	registro y consulta de NNA's, gestión de expedientes, registro de 
	valoraciones y consulta del calendario de casos.
\end{description}

Ambos flujos comparten elementos comunes de navegación presentes en todas 
las pantallas del sistema:

\begin{description}
	\item[Barra superior:] Presente en todas las pantallas posteriores al 
	inicio de sesión. Muestra el logotipo del sistema SICORRE, el nombre 
	completo del usuario autenticado, su rol y el botón 
	\IUbutton{Cerrar sesión}.
\end{description}

La navegación entre pantallas se muestra en la figura~\ref{fig:mapa}, en 
ella se ilustra el flujo completo de interacción del usuario con el sistema, 
desde el acceso inicial hasta las funcionalidades de cada módulo, 
diferenciando los caminos disponibles según el rol autenticado.

\begin{figure}[htbp!]
	\begin{center}
		\includegraphics[width=.95\textwidth]{images/flujo}
		\caption{Mapa de navegación del sistema SICORRE.}
		\label{fig:mapa}
	\end{center}
\end{figure}

\clearpage
%---------------------------------------------------------


%--------------------------------------
\section{IU01 Pantalla de Inicio de Sesión}

\subsection{Objetivo}
Permitir al personal autorizado de la Fundación Futuro con Derecho 
autenticarse en el sistema SICORRE mediante sus credenciales de acceso 
(correo electrónico y contraseña), redirigiendo al usuario hacia el 
panel correspondiente a su rol una vez validada la identidad.

\subsection{Diseño}
Esta pantalla \IUref{IU01}{Pantalla de Inicio de Sesión} es la primera 
que aparece al acceder al sistema SICORRE. El usuario debe ingresar su 
correo electrónico institucional y su contraseña en los campos 
correspondientes y presionar el botón \IUbutton{Ingresar} para 
autenticarse. En caso de que las credenciales sean incorrectas, el 
sistema mostrará un mensaje de error sin especificar cuál de los dos 
campos es incorrecto, por razones de seguridad. Si la autenticación 
es exitosa, el sistema redirigirá automáticamente al usuario hacia 
su panel según su rol: la directora y coordinadora serán dirigidas 
a la \IUref{IU02}{Pantalla del Panel de Administrador}, mientras que 
los especialistas serán dirigidos a la 
\IUref{IU03}{Pantalla del Panel de Especialista}.

\IUfig[.7]{IU1-IniciarSesion}{IU01}{Pantalla de Inicio de Sesión.}

\subsection{Salidas}
\begin{itemize}
	\item Mensaje de bienvenida con el nombre del usuario autenticado.
	\item Redirección a la \IUref{IU02}{Pantalla del Panel de 
		Administrador} si el usuario tiene rol de administrador 
	(directora o coordinadora).
	\item Redirección a la Pantalla del Panel de 
		Especialista (por programar) si el usuario tiene rol de especialista 
	(trabajador social, abogado, psicólogo o médico).
	\item Mensaje de error \MSGref{MSG01}{Credenciales incorrectas} 
	si el correo o la contraseña no son válidos.
	\item Mensaje de error \MSGref{MSG02}{Usuario inactivo} si la 
	cuenta del usuario ha sido revocada por el administrador.
\end{itemize}

\subsection{Entradas}
\begin{itemize}
	\item \textbf{Correo electrónico:} Campo de texto donde el usuario 
	ingresa su dirección de correo electrónico institucional registrada 
	en el sistema.
	\item \textbf{Contraseña:} Campo de texto enmascarado donde el 
	usuario ingresa su contraseña de acceso.
\end{itemize}

\subsection{Comandos}
\begin{itemize}
	\item \IUbutton{Ingresar}: Valida las credenciales ingresadas 
	contra la base de datos del sistema. Si la validación es exitosa 
	y el usuario tiene rol de administrador, muestra la 
	\IUref{IU02}{Pantalla del Panel de Administrador}. Si el usuario 
	tiene rol de especialista, muestra la 
	Pantalla del Panel de Especialista. Si las 
	credenciales son incorrectas, muestra el mensaje 
	\MSGref{MSG01}{Credenciales incorrectas}. Si la cuenta está 
	revocada, muestra el mensaje \MSGref{MSG02}{Usuario inactivo}.
	
	\item \IUbutton{Cancelar}: Limpia los campos de correo electrónico 
	y contraseña, permitiendo al usuario reingresar sus datos.
\end{itemize}


\section{IU02 Pantalla del Panel de Administrador}

\subsection{Objetivo}
Proporcionar a la directora y coordinadora un panel central de navegación 
desde el cual puedan acceder a los cuatro módulos principales del sistema 
SICORRE disponibles para el rol de administrador.

\subsection{Diseño}
Esta pantalla \IUref{IU02}{Pantalla del Panel de Administrador} aparece 
inmediatamente después de que la directora o coordinadora inicia sesión 
correctamente desde la \IUref{IU01}{Pantalla de Inicio de Sesión}. 

En la parte superior se muestra la barra de navegación con el nombre del 
usuario autenticado, su rol y el botón \IUbutton{Cerrar sesión}. El área 
central presenta cuatro botones de acceso directo a los módulos del sistema, 
organizados de forma clara para facilitar la navegación sin necesidad de 
conocimientos técnicos previos.

\IUfig[.9]{DashAdmin}{IU02}{Pantalla del Panel de Administrador.}

\subsection{Salidas}
\begin{itemize}
	\item Nombre y rol del usuario autenticado en la barra superior.
\end{itemize}

\subsection{Entradas}
\begin{itemize}
	\item No requiere entradas del usuario. La información del usuario 
	autenticado se carga automáticamente al iniciar sesión.
\end{itemize}

\subsection{Comandos}
\begin{itemize}
	\item \IUbutton{Consultar usuarios}: Redirige a la 
	\IUref{IU03}{Pantalla de Consulta de Usuarios} donde la directora 
	o coordinadora puede visualizar, modificar, revocar y eliminar 
	el acceso del personal registrado en el sistema.
	
	\item \IUbutton{Consultar equipos}: Redirige a la 
	\IUref{IU04}{Pantalla de Consulta de Equipos Multidisciplinarios} 
	donde se pueden visualizar y gestionar los equipos de trabajo 
	registrados en el sistema.
	
	\item \IUbutton{Consultar NNAs y Tutores}: Redirige a la 
	\IUref{IU05}{Pantalla de Consulta de NNAs y Tutores} donde se 
	puede visualizar el listado de niñas, niños y adolescentes 
	registrados en el sistema junto con los datos de sus tutores.
	
	\item \IUbutton{Restitución de Derechos}: Redirige a la 
	\IUref{IU06}{Pantalla de Restitución de Derechos} donde se 
	puede consultar el estado de los expedientes activos y los 
	derechos vulnerados identificados para cada NNA.
	
	\item \IUbutton{Cerrar sesión}: Cierra la sesión activa del 
	usuario y regresa a la \IUref{IU01}{Pantalla de Inicio de Sesión}.
\end{itemize}


% !TeX root = ../proyecto.tex
%--------------------------------------
\section{IU03 Pantalla de Consulta de Usuarios}

\subsection{Objetivo}
Permitir a la directora y coordinadora visualizar el listado completo 
del personal registrado en el sistema SICORRE, así como acceder a las 
opciones de gestión de cada usuario: modificar sus datos, revocar su 
acceso o eliminar su registro. También permite iniciar el flujo de 
registro de un nuevo integrante del equipo.

\subsection{Diseño}
Esta pantalla \IUref{IU03}{Pantalla de Consulta de Usuarios} es 
accesible desde el botón \IUbutton{Consultar usuarios} de la 
\IUref{IU02}{Pantalla del Panel de Administrador}. 

Presenta una tabla con el listado de todo el personal registrado 
en el sistema. Cada fila corresponde a un usuario e incluye su 
nombre completo, rol asignado dentro del equipo multidisciplinario, 
estatus de cuenta y las opciones de acción disponibles. En la parte 
superior de la pantalla se encuentra el botón 
\IUbutton{Registrar personal} para dar de alta un nuevo integrante.

\IUfig[.9]{ListaUsuarios}{IU03}{Pantalla de Consulta de Usuarios.}

\subsection{Salidas}
\begin{itemize}
	\item Tabla con el listado del personal registrado en el sistema, 
	mostrando por cada usuario:
	\begin{itemize}
		\item \textbf{Nombre:} Nombre completo del usuario registrado.
		\item \textbf{Rol:} Perfil asignado al usuario dentro del 
		equipo multidisciplinario (Administrador, Trabajador Social, 
		Abogado, Psicólogo o Médico).
		\item \textbf{Activo:} Estatus actual de la cuenta del usuario 
		(Activo o Revocado).
		\item \textbf{Opciones:} Acciones disponibles para cada usuario 
		(Modificar, Revocar, Eliminar).
	\end{itemize}
	\item Mensaje \MSGref{MSG03}{Sin usuarios registrados} en caso de 
	que no exista personal registrado en el sistema.
\end{itemize}

\subsection{Entradas}
\begin{itemize}
	\item No requiere entradas directas del usuario. El listado se 
	carga automáticamente al acceder a la pantalla.
\end{itemize}

\subsection{Comandos}
\begin{itemize}
	\item \IUbutton{Registrar personal}: Redirige a la 
	\IUref{IU04}{Pantalla de Registro de Personal} donde el 
	administrador puede capturar los datos de un nuevo integrante 
	del equipo multidisciplinario.
	
	\item \IUbutton{Modificar}: Disponible por cada fila de usuario 
	en la tabla. Redirige a la 
	\IUref{IU05}{Pantalla de Modificación de Usuario} donde se 
	pueden editar los datos del usuario seleccionado.
	
	\item \IUbutton{Revocar}: Disponible por cada fila de usuario 
	activo en la tabla. Muestra un mensaje de confirmación 
	\MSGref{MSG04}{Confirmar revocación de acceso} antes de 
	proceder. Al confirmar, el estatus del usuario cambia a 
	\textit{Revocado} y se actualiza en la tabla sin eliminar 
	su historial de actividad en el sistema.
	
	\item \IUbutton{Eliminar}: Disponible por cada fila de usuario 
	en la tabla. Muestra un mensaje de confirmación 
	\MSGref{MSG05}{Confirmar eliminación de usuario} antes de 
	proceder. Al confirmar, el registro del usuario es eliminado 
	permanentemente del sistema.
	
	\item \IUbutton{Volver}: Regresa a la 
	\IUref{IU02}{Pantalla del Panel de Administrador}.
\end{itemize}


% !TeX root = ../proyecto.tex
%--------------------------------------
\section{IU04 Pantalla de Registro de Personal}

\subsection{Objetivo}
Permitir a la directora y coordinadora registrar un nuevo integrante 
del equipo multidisciplinario en el sistema SICORRE, capturando sus 
datos personales, de domicilio y de acceso al sistema.

\subsection{Diseño}
Esta pantalla \IUref{IU04}{Pantalla de Registro de Personal} es 
accesible desde el botón \IUbutton{Registrar personal} de la 
\IUref{IU03}{Pantalla de Consulta de Usuarios}.

El formulario se organiza en dos secciones principales: 
\textbf{Datos personales}, que incluye nombre, CURP, RFC, fecha 
de nacimiento, sexo, tipo de personal, empleado, estado de cuenta 
y rol; y \textbf{Dirección}, que utiliza el catálogo de códigos 
postales de SEPOMEX para autocompletar el estado, municipio y 
colonia una vez ingresado el código postal. Al presionar 
\IUbutton{Buscar} con un código postal válido, los campos de 
Estado y Municipio/Alcaldía se llenan automáticamente y el campo 
Colonia muestra un selector con las colonias disponibles para 
dicho código postal. Finalmente se capturan los datos de acceso 
al sistema: correo electrónico y contraseña provisional.

\IUfig[.9]{RegistroPersonal}{IU04}{Pantalla de Registro de Personal.}

\subsection{Salidas}
\begin{itemize}
	\item Campos de \textbf{Estado} y \textbf{Municipio/Alcaldía} 
	autocompletados al ingresar un código postal válido y presionar 
	\IUbutton{Buscar}.
	\item Selector de \textbf{Colonia} poblado con las colonias 
	correspondientes al código postal ingresado.
	\item Mensaje \MSGref{MSG06}{Registro exitoso} al guardar 
	correctamente el nuevo usuario.
	\item Mensaje \MSGref{MSG07}{CURP ya registrada} si la CURP 
	ingresada ya existe en el sistema.
	\item Mensaje \MSGref{MSG08}{Correo ya registrado} si el correo 
	electrónico ingresado ya está asociado a otro usuario.
	\item Mensaje \MSGref{MSG09}{Código postal no encontrado} si el 
	código postal ingresado no existe en el catálogo SEPOMEX.
	\item Mensaje \MSGref{MSG10}{Campos obligatorios incompletos} si 
	el usuario intenta guardar sin completar todos los campos 
	requeridos.
\end{itemize}

\subsection{Entradas}
\begin{itemize}
	\item \textbf{Nombre(s):} Campo de texto obligatorio. Nombre o 
	nombres del nuevo integrante.
	\item \textbf{Apellido Paterno:} Campo de texto obligatorio.
	\item \textbf{Apellido Materno:} Campo de texto obligatorio.
	\item \textbf{CURP:} Campo de texto obligatorio. Debe tener 
	18 caracteres con formato válido de CURP mexicana.
	\item \textbf{RFC:} Campo de texto obligatorio. Debe tener 
	13 caracteres con formato válido de RFC mexicano.
	\item \textbf{Fecha de nacimiento:} Campo de fecha obligatorio 
	con formato dd/mm/aaaa.
	\item \textbf{Seleccionar Sexo:} Lista desplegable obligatoria 
	con las opciones disponibles.
	\item \textbf{Tipo de personal:} Lista desplegable obligatoria. 
	Por defecto muestra la opción \textit{Empleado}.
	\item \textbf{Estado de cuenta:} Lista desplegable obligatoria. 
	Por defecto muestra la opción \textit{Activo}.
	\item \textbf{Calle:} Campo de texto obligatorio correspondiente 
	a la calle del domicilio.
	\item \textbf{Número Exterior:} Campo de texto obligatorio.
	\item \textbf{Número Interior:} Campo de texto opcional.
	\item \textbf{Código Postal:} Campo de texto obligatorio. 
	Acepta 5 dígitos numéricos. Ejemplo: 06600.
	\item \textbf{Estado:} Campo de texto de solo lectura. 
	Se autocompleta al buscar el código postal.
	\item \textbf{Municipio / Alcaldía:} Campo de texto de solo 
	lectura. Se autocompleta al buscar el código postal.
	\item \textbf{Colonia:} Lista desplegable que se habilita 
	únicamente después de realizar una búsqueda de código postal 
	exitosa. Muestra el mensaje \textit{Busca un código postal 
		primero} mientras no se haya realizado la búsqueda.
	\item \textbf{Tipo de Vivienda:} Lista desplegable obligatoria 
	con los tipos de vivienda disponibles en el catálogo del sistema.
	\item \textbf{Rol:} Lista desplegable obligatoria con los roles 
	disponibles en el sistema (Administrador, Trabajador Social, 
	Abogado, Psicólogo, Médico).
	\item \textbf{Correo:} Campo de texto obligatorio con formato 
	de correo electrónico válido.
	\item \textbf{Contraseña Provisional:} Campo de texto enmascarado 
	obligatorio que será asignada al nuevo usuario para su primer 
	acceso al sistema.
\end{itemize}

\subsection{Comandos}
\begin{itemize}
	\item \IUbutton{Buscar}: Consulta el catálogo SEPOMEX con el 
	código postal ingresado. Si es válido, autocompleta los campos 
	de Estado y Municipio/Alcaldía y habilita el selector de Colonia 
	con las opciones correspondientes. Si no es válido, muestra el 
	mensaje \MSGref{MSG09}{Código postal no encontrado}.
	
	\item \IUbutton{Guardar}: Valida que todos los campos obligatorios 
	estén completos y con formato correcto. Si la validación es 
	exitosa, registra al nuevo usuario en el sistema, muestra el 
	mensaje \MSGref{MSG06}{Registro exitoso} y regresa a la 
	\IUref{IU03}{Pantalla de Consulta de Usuarios}. Si hay campos 
	incompletos muestra el mensaje 
	\MSGref{MSG10}{Campos obligatorios incompletos}. Si la CURP ya 
	existe muestra \MSGref{MSG07}{CURP ya registrada}. Si el correo 
	ya existe muestra \MSGref{MSG08}{Correo ya registrado}.
	
	\item \IUbutton{Cancelar}: Descarta todos los datos ingresados 
	en el formulario y regresa a la 
	\IUref{IU03}{Pantalla de Consulta de Usuarios} sin realizar 
	ningún cambio en el sistema.
\end{itemize}


% !TeX root = ../proyecto.tex
%--------------------------------------
\section{IU05 Pantalla de Modificación de Usuario}

\subsection{Objetivo}
Permitir a la directora y coordinadora editar los datos personales, 
de domicilio y de acceso de un integrante del equipo multidisciplinario 
previamente registrado en el sistema SICORRE.

\subsection{Diseño}
Esta pantalla \IUref{IU05}{Pantalla de Modificación de Usuario} es 
accesible desde el botón \IUbutton{Modificar} correspondiente a un 
usuario en la \IUref{IU03}{Pantalla de Consulta de Usuarios}.

Al acceder, el sistema recupera automáticamente los datos actuales 
del usuario seleccionado y los presenta precargados en todos los 
campos del formulario, permitiendo al administrador identificar 
y modificar únicamente la información que requiera actualización. 

La estructura del formulario es idéntica a la de la 
\IUref{IU04}{Pantalla de Registro de Personal}, organizada en dos 
secciones: \textbf{Datos personales} y \textbf{Dirección}. El campo 
de Colonia se muestra ya seleccionado con el valor actual del usuario, 
sin necesidad de volver a buscar el código postal a menos que se 
desee cambiar el domicilio. El campo de contraseña se muestra vacío 
por seguridad; si se deja vacío al guardar, la contraseña actual 
del usuario se conserva sin cambios.

\IUfig[.9]{ModificarUsuario}{IU05}{Pantalla de Modificación de Usuario.}

\subsection{Salidas}
\begin{itemize}
	\item Todos los campos del formulario precargados con los datos 
	actuales del usuario seleccionado.
	\item Campos de \textbf{Estado} y \textbf{Municipio/Alcaldía} 
	autocompletados si se realiza una nueva búsqueda por código postal.
	\item Selector de \textbf{Colonia} actualizado si se realiza 
	una nueva búsqueda por código postal.
	\item Mensaje \MSGref{MSG11}{Modificación exitosa} al guardar 
	los cambios correctamente.
	\item Mensaje \MSGref{MSG07}{CURP ya registrada} si la CURP 
	modificada ya pertenece a otro usuario en el sistema.
	\item Mensaje \MSGref{MSG08}{Correo ya registrado} si el correo 
	modificado ya está asociado a otro usuario en el sistema.
	\item Mensaje \MSGref{MSG09}{Código postal no encontrado} si el 
	código postal ingresado no existe en el catálogo SEPOMEX.
	\item Mensaje \MSGref{MSG10}{Campos obligatorios incompletos} si 
	se intenta guardar sin completar todos los campos requeridos.
\end{itemize}

\subsection{Entradas}
\begin{itemize}
	\item \textbf{Nombre(s):} Campo de texto obligatorio. 
	Precargado con el valor actual.
	\item \textbf{Apellido Paterno:} Campo de texto obligatorio. 
	Precargado con el valor actual.
	\item \textbf{Apellido Materno:} Campo de texto obligatorio. 
	Precargado con el valor actual.
	\item \textbf{CURP:} Campo de texto obligatorio de 18 caracteres. 
	Precargado con el valor actual.
	\item \textbf{RFC:} Campo de texto obligatorio de 13 caracteres. 
	Precargado con el valor actual.
	\item \textbf{Fecha de nacimiento:} Campo de fecha obligatorio 
	con formato dd/mm/aaaa. Precargado con el valor actual.
	\item \textbf{Seleccionar Sexo:} Lista desplegable obligatoria. 
	Precargada con el valor actual.
	\item \textbf{Tipo de personal:} Lista desplegable obligatoria. 
	Precargada con el valor actual.
	\item \textbf{Estado de cuenta:} Lista desplegable obligatoria. 
	Precargada con el valor actual.
	\item \textbf{Calle:} Campo de texto obligatorio. 
	Precargado con el valor actual.
	\item \textbf{Número Exterior:} Campo de texto obligatorio. 
	Precargado con el valor actual.
	\item \textbf{Número Interior:} Campo de texto opcional. 
	Precargado con el valor actual si existe.
	\item \textbf{Código Postal:} Campo de texto obligatorio de 
	5 dígitos. Precargado con el valor actual.
	\item \textbf{Estado:} Campo de solo lectura. Precargado con 
	el valor actual. Se actualiza si se realiza una nueva búsqueda 
	de código postal.
	\item \textbf{Municipio / Alcaldía:} Campo de solo lectura. 
	Precargado con el valor actual. Se actualiza si se realiza 
	una nueva búsqueda de código postal.
	\item \textbf{Colonia:} Lista desplegable obligatoria. 
	Precargada con el valor actual del usuario.
	\item \textbf{Tipo de Vivienda:} Lista desplegable obligatoria. 
	Precargada con el valor actual.
	\item \textbf{Rol:} Lista desplegable obligatoria. 
	Precargada con el valor actual.
	\item \textbf{Correo:} Campo de texto obligatorio. 
	Precargado con el valor actual.
	\item \textbf{Contraseña:} Campo de texto enmascarado opcional. 
	Se muestra vacío por seguridad. Si se deja vacío al guardar, 
	la contraseña actual del usuario se conserva sin cambios.
\end{itemize}

\subsection{Comandos}
\begin{itemize}
	\item \IUbutton{Buscar}: Consulta el catálogo SEPOMEX con el 
	código postal ingresado. Si es válido, actualiza los campos de 
	Estado y Municipio/Alcaldía y recarga el selector de Colonia 
	con las opciones correspondientes. Si no es válido, muestra 
	el mensaje \MSGref{MSG09}{Código postal no encontrado}.
	
	\item \IUbutton{Guardar cambios}: Valida que todos los campos 
	obligatorios estén completos y con formato correcto. Si la 
	validación es exitosa, actualiza los datos del usuario en el 
	sistema, muestra el mensaje \MSGref{MSG11}{Modificación exitosa} 
	y regresa a la \IUref{IU03}{Pantalla de Consulta de Usuarios}. 
	Si hay campos incompletos muestra 
	\MSGref{MSG10}{Campos obligatorios incompletos}. Si la CURP 
	modificada pertenece a otro usuario muestra 
	\MSGref{MSG07}{CURP ya registrada}. Si el correo modificado 
	pertenece a otro usuario muestra 
	\MSGref{MSG08}{Correo ya registrado}.
	
	\item \IUbutton{Cancelar}: Descarta todos los cambios realizados 
	en el formulario y regresa a la 
	\IUref{IU03}{Pantalla de Consulta de Usuarios} conservando los 
	datos originales del usuario sin modificación alguna.
\end{itemize}

% !TeX root = ../proyecto.tex
%--------------------------------------
\section{IU06 Pantalla de Consulta de Usuario}

\subsection{Objetivo}
Permitir a la directora y coordinadora visualizar el detalle completo 
de los datos de un integrante del equipo multidisciplinario registrado 
en el sistema SICORRE, así como acceder a las opciones de gestión 
disponibles para dicho usuario.

\subsection{Diseño}
Esta pantalla \IUref{IU06}{Pantalla de Consulta de Usuario} es 
accesible al seleccionar el nombre de un usuario en la tabla de la 
\IUref{IU03}{Pantalla de Consulta de Usuarios}.

Al acceder, el sistema recupera y presenta todos los datos del usuario 
seleccionado organizados en tres secciones de solo lectura:

\begin{description}
	\item[Datos personales:] Nombre completo, CURP, RFC, fecha de 
	nacimiento, sexo, tipo de personal y estado de cuenta.
	\item[Dirección:] Calle, número exterior, número interior, 
	código postal, estado, municipio o alcaldía, colonia y tipo 
	de vivienda.
	\item[Datos del sistema:] Rol asignado, correo electrónico 
	y fecha de registro en el sistema.
\end{description}

En la parte inferior de la pantalla se presentan los botones de 
acción disponibles para gestionar al usuario consultado.

\IUfig[.9]{ConsultaUsaurios.png}{IU06}{Pantalla de Consulta de Usuario.}

\subsection{Salidas}
\begin{itemize}
	\item Datos personales completos del usuario seleccionado 
	en modo de solo lectura.
	\item Datos de domicilio completos del usuario seleccionado 
	en modo de solo lectura.
	\item Datos del sistema del usuario seleccionado en modo 
	de solo lectura.
	\item Mensaje \MSGref{MSG04}{Confirmar revocación de acceso} 
	al presionar \IUbutton{Revocar}.
	\item Mensaje \MSGref{MSG05}{Confirmar eliminación de usuario} 
	al presionar \IUbutton{Eliminar}.
\end{itemize}

\subsection{Entradas}
\begin{itemize}
	\item No requiere entradas del usuario. Todos los datos se 
	cargan automáticamente a partir del usuario seleccionado 
	en la \IUref{IU03}{Pantalla de Consulta de Usuarios}.
\end{itemize}

\subsection{Comandos}
\begin{itemize}
	\item \IUbutton{Modificar}: Redirige a la 
	\IUref{IU05}{Pantalla de Modificación de Usuario} con los 
	datos del usuario actual precargados en el formulario.
	
	\item \IUbutton{Revocar}: Muestra el mensaje de confirmación 
	\MSGref{MSG04}{Confirmar revocación de acceso}. Al confirmar, 
	el estatus de la cuenta del usuario cambia a \textit{Revocado}, 
	impidiendo su acceso al sistema sin eliminar su historial de 
	actividad. La pantalla se actualiza reflejando el nuevo estatus.
	
	\item \IUbutton{Eliminar}: Muestra el mensaje de confirmación 
	\MSGref{MSG05}{Confirmar eliminación de usuario}. Al confirmar, 
	el registro del usuario es eliminado permanentemente del sistema 
	y redirige a la 
	\IUref{IU03}{Pantalla de Consulta de Usuarios}.
	
	\item \IUbutton{Volver}: Regresa a la 
	\IUref{IU03}{Pantalla de Consulta de Usuarios} sin realizar 
	ningún cambio.
\end{itemize}


% !TeX root = ../proyecto.tex
%--------------------------------------
\section{IU07 Pantalla de Consulta de Equipos Multidisciplinarios}

\subsection{Objetivo}
Permitir a la directora y coordinadora visualizar el listado completo 
de los equipos multidisciplinarios registrados en el sistema SICORRE, 
así como acceder a las opciones de gestión de cada equipo.

\subsection{Diseño}
Esta pantalla \IUref{IU07}{Pantalla de Consulta de Equipos 
	Multidisciplinarios} es accesible desde el botón 
\IUbutton{Consultar equipos} de la 
\IUref{IU02}{Pantalla del Panel de Administrador}.

Presenta una tabla con el listado de todos los equipos 
multidisciplinarios registrados en el sistema. Cada fila corresponde 
a un equipo e incluye su nombre, estatus, el listado de integrantes 
que lo conforman y las opciones de acción disponibles. En la parte 
superior de la pantalla se encuentra el botón 
\IUbutton{Registrar equipo} para dar de alta un nuevo equipo 
multidisciplinario.

\IUfig[.9]{ConsultaEquipo}{IU07}{Pantalla de Consulta de Equipos 
	Multidisciplinarios.}

\subsection{Salidas}
\begin{itemize}
	\item Tabla con el listado de equipos multidisciplinarios 
	registrados en el sistema, mostrando por cada equipo:
	\begin{itemize}
		\item \textbf{Nombre:} Nombre del equipo multidisciplinario.
		\item \textbf{Estatus:} Estado actual del equipo 
		(Activo o Inactivo).
		\item \textbf{Integrantes:} Listado de los integrantes 
		que conforman el equipo con su rol correspondiente.
		\item \textbf{Opciones:} Acciones disponibles para cada 
		equipo (Modificar, Eliminar).
	\end{itemize}
	\item Mensaje \MSGref{MSG12}{Sin equipos registrados} en caso 
	de que no exista ningún equipo registrado en el sistema.
\end{itemize}

\subsection{Entradas}
\begin{itemize}
	\item No requiere entradas directas del usuario. El listado 
	se carga automáticamente al acceder a la pantalla.
\end{itemize}

\subsection{Comandos}
\begin{itemize}
	\item \IUbutton{Registrar equipo}: Redirige a la 
	\IUref{IU08}{Pantalla de Registro de Equipo} donde el 
	administrador puede capturar los datos de un nuevo equipo 
	multidisciplinario y asignar sus integrantes.
	
	\item \IUbutton{Modificar}: Disponible por cada fila de equipo 
	en la tabla. Redirige a la 
	\IUref{IU09}{Pantalla de Modificación de Equipo} donde se 
	pueden editar los datos del equipo seleccionado y gestionar 
	sus integrantes.
	
	\item \IUbutton{Eliminar}: Disponible por cada fila de equipo 
	en la tabla. Muestra el mensaje de confirmación 
	\MSGref{MSG13}{Confirmar eliminación de equipo} antes de 
	proceder. Al confirmar, el registro del equipo es eliminado 
	permanentemente del sistema y la tabla se actualiza.
	
	\item \IUbutton{Volver}: Regresa a la 
	\IUref{IU02}{Pantalla del Panel de Administrador}.
\end{itemize}



% !TeX root = ../proyecto.tex
%--------------------------------------
\section{IU08 Pantalla de Registro de Equipo}

\subsection{Objetivo}
Permitir a la directora y coordinadora registrar un nuevo equipo 
multidisciplinario en el sistema SICORRE asignándole un nombre 
único que lo identifique.

\subsection{Diseño}
Esta pantalla \IUref{IU08}{Pantalla de Registro de Equipo} es 
accesible desde el botón \IUbutton{Registrar equipo} de la 
\IUref{IU07}{Pantalla de Consulta de Equipos Multidisciplinarios}.

Presenta un formulario simple con un único campo de texto para 
capturar el nombre del nuevo equipo multidisciplinario y los 
botones de acción correspondientes.

\IUfig[.9]{RegistroEquipo}{IU08}{Pantalla de Registro de Equipo.}

\subsection{Salidas}
\begin{itemize}
	\item Mensaje \MSGref{MSG14}{Equipo registrado exitosamente} 
	al guardar correctamente el nuevo equipo.
	\item Mensaje \MSGref{MSG15}{Nombre de equipo ya registrado} 
	si el nombre ingresado ya existe en el sistema.
	\item Mensaje \MSGref{MSG10}{Campos obligatorios incompletos} 
	si se intenta guardar sin ingresar el nombre del equipo.
\end{itemize}

\subsection{Entradas}
\begin{itemize}
	\item \textbf{Nombre del equipo:} Campo de texto obligatorio 
	con el nombre que identificará al nuevo equipo 
	multidisciplinario dentro del sistema.
\end{itemize}

\subsection{Comandos}
\begin{itemize}
	\item \IUbutton{Guardar}: Valida que el campo de nombre esté 
	completo y no esté duplicado. Si la validación es exitosa, 
	registra el nuevo equipo en el sistema con estatus 
	\textit{Activo}, muestra el mensaje 
	\MSGref{MSG14}{Equipo registrado exitosamente} y regresa a 
	la \IUref{IU07}{Pantalla de Consulta de Equipos 
		Multidisciplinarios}. Si el campo está vacío muestra 
	\MSGref{MSG10}{Campos obligatorios incompletos}. Si el nombre 
	ya existe muestra 
	\MSGref{MSG15}{Nombre de equipo ya registrado}.
	
	\item \IUbutton{Cancelar}: Descarta el dato ingresado y 
	regresa a la \IUref{IU07}{Pantalla de Consulta de Equipos 
		Multidisciplinarios} sin realizar ningún cambio en el sistema.
\end{itemize}


% !TeX root = ../proyecto.tex
%--------------------------------------
\section{IU09 Pantalla de Añadir Integrantes a Equipo}

\subsection{Objetivo}
Permitir a la directora y coordinadora asignar integrantes 
disponibles a un equipo multidisciplinario registrado en el 
sistema SICORRE, mostrando únicamente el personal que no 
pertenece a ningún equipo activo.

\subsection{Diseño}
Esta pantalla \IUref{IU09}{Pantalla de Añadir Integrantes a Equipo} 
es accesible desde el botón \IUbutton{Modificar} de un equipo en la 
\IUref{IU07}{Pantalla de Consulta de Equipos Multidisciplinarios}.

Presenta una tabla con el listado del personal disponible, es decir, 
los integrantes registrados en el sistema que actualmente no 
pertenecen a ningún equipo multidisciplinario activo. Cada fila 
muestra el nombre completo del integrante y su rol, junto con un 
botón \IUbutton{Agregar} para asignarlo al equipo actual.

En caso de que todo el personal ya esté asignado a un equipo, 
la tabla mostrará el mensaje 
\MSGref{MSG16}{No hay integrantes disponibles}.

\IUfig[.9]{AnadirIntegrantes}{IU09}{Pantalla de Añadir Integrantes 
	a Equipo.}

\subsection{Salidas}
\begin{itemize}
	\item Nombre del equipo al que se están asignando integrantes, 
	mostrado en el encabezado de la pantalla.
	\item Tabla con el listado del personal disponible mostrando 
	por cada integrante:
	\begin{itemize}
		\item \textbf{Nombre:} Nombre completo del integrante 
		disponible.
		\item \textbf{Rol:} Perfil del integrante dentro del 
		equipo multidisciplinario.
		\item \textbf{Opciones:} Botón \IUbutton{Agregar} 
		por cada fila.
	\end{itemize}
	\item Mensaje \MSGref{MSG16}{No hay integrantes disponibles} 
	si todo el personal activo ya pertenece a un equipo.
	\item Mensaje \MSGref{MSG17}{Integrante agregado exitosamente} 
	al asignar correctamente un integrante al equipo.
\end{itemize}

\subsection{Entradas}
\begin{itemize}
	\item No requiere entradas directas del usuario. El listado 
	de integrantes disponibles se carga automáticamente al 
	acceder a la pantalla.
\end{itemize}

\subsection{Comandos}
\begin{itemize}
	\item \IUbutton{Agregar}: Disponible por cada fila de 
	integrante en la tabla. Asigna al integrante seleccionado 
	al equipo actual, muestra el mensaje 
	\MSGref{MSG17}{Integrante agregado exitosamente} y actualiza 
	la tabla eliminando al integrante recién asignado de la 
	lista de disponibles.
	
	\item \IUbutton{Volver}: Regresa a la 
	\IUref{IU07}{Pantalla de Consulta de Equipos 
		Multidisciplinarios} sin realizar cambios adicionales.
\end{itemize}


% !TeX root = ../proyecto.tex
%--------------------------------------\section{IU10 Pantalla de Modificación de Equipo}

\subsection{Objetivo}
Permitir a la directora y coordinadora modificar el nombre de un equipo multidisciplinario registrado en el sistema SICORRE, así como eliminar integrantes que actualmente pertenecen a dicho equipo.

\subsection{Diseño}
Esta pantalla \IUref{IU10}{Pantalla de Modificación de Equipo} es accesible desde el botón \IUbutton{Modificar} de un equipo en la \IUref{IU07}{Pantalla de Consulta de Equipos Multidisciplinarios}. Presenta un campo de texto editable con el nombre actual del equipo, permitiendo al usuario actualizarlo. Además, muestra una tabla con los integrantes que actualmente pertenecen al equipo, donde cada fila contiene el nombre completo del integrante, su rol y un botón \IUbutton{Eliminar} para removerlo del equipo. En caso de que el equipo no cuente con integrantes registrados, la tabla mostrará el mensaje \MSGref{MSG18}{No hay integrantes en el equipo}.

\IUfig[.9]{ModificarEquipo}{IU10}{Pantalla de Modificación de Equipo.}

\subsection{Salidas}
\begin{itemize}
	\item Campo de texto con el nombre actual del equipo, disponible para edición.
	\item Tabla con el listado de integrantes actuales del equipo, mostrando por cada integrante:
	\begin{itemize}
		\item \textbf{Nombre:} Nombre completo del integrante.
		\item \textbf{Rol:} Perfil del integrante dentro del equipo multidisciplinario.
		\item \textbf{Opciones:} Botón \IUbutton{Eliminar} por cada fila.
	\end{itemize}
	\item Mensaje \MSGref{MSG18}{No hay integrantes en el equipo} si el equipo no tiene integrantes asignados.
	\item Mensaje \MSGref{MSG19}{Equipo modificado exitosamente} al guardar los cambios correctamente.
	\item Mensaje \MSGref{MSG20}{Integrante eliminado exitosamente} al remover correctamente un integrante del equipo.
\end{itemize}

\subsection{Entradas}
\begin{itemize}
	\item \textbf{Nombre del equipo:} Campo de texto editable con el nombre actual del equipo. El usuario puede modificarlo para actualizarlo.
\end{itemize}

\subsection{Comandos}
\begin{itemize}
	\item \IUbutton{Guardar}: Almacena el nuevo nombre del equipo en el sistema, muestra el mensaje \MSGref{MSG19}{Equipo modificado exitosamente} y permanece en la pantalla actual con los datos actualizados.
	\item \IUbutton{Eliminar}: Disponible por cada fila de integrante en la tabla. Remueve al integrante seleccionado del equipo, muestra el mensaje \MSGref{MSG20}{Integrante eliminado exitosamente} y actualiza la tabla eliminando la fila correspondiente.
	\item \IUbutton{Volver}: Regresa a la \IUref{IU07}{Pantalla de Consulta de Equipos Multidisciplinarios} sin realizar cambios adicionales.
\end{itemize}


% !TeX root = ../proyecto.tex
%--------------------------------------
\section{IU11 Pantalla de Consulta de NNAs y Tutores}

\subsection{Objetivo}
Permitir a la directora y coordinadora consultar el listado de Niñas, Niños y Adolescentes (NNAs) y sus tutores registrados en el sistema SICORRE, así como acceder a la edición o eliminación de los registros de NNAs.

\subsection{Diseño}
Esta pantalla \IUref{IU11}{Pantalla de Consulta de NNAs y Tutores} es accesible desde la \IUref{IU02}{Pantalla del Panel de Administrador}. Presenta dos tablas independientes: la primera corresponde al listado de NNAs y la segunda al listado de tutores registrados en el sistema.

La tabla de NNAs muestra por cada registro el nombre completo, CURP, fecha de nacimiento, sexo, estatus escolar, equipo asignado y estatus, junto con dos columnas de operaciones: \textbf{Acciones}, que contiene el botón \IUbutton{Editar}, y \textbf{Gestión}, que contiene el botón \IUbutton{Eliminar}.

La tabla de Tutores muestra por cada registro el nombre completo, CURP, parentesco y sexo, siendo únicamente de consulta sin acciones disponibles.

En caso de no existir NNAs registrados, la tabla correspondiente mostrará el mensaje \MSGref{MSG21}{No hay NNAs registrados}. De igual forma, si no existen tutores registrados, se mostrará el mensaje \MSGref{MSG22}{No hay tutores registrados}.

\IUfig[.9]{ConsultaNNAsTutores}{IU11}{Pantalla de Consulta de NNAs y Tutores.}

\subsection{Salidas}
\begin{itemize}
	\item Tabla de NNAs con los siguientes campos por registro:
	\begin{itemize}
		\item \textbf{Nombre:} Nombre completo del NNA.
		\item \textbf{CURP:} Clave Única de Registro de Población del NNA.
		\item \textbf{Fecha de nacimiento:} Fecha de nacimiento del NNA.
		\item \textbf{Sexo:} Sexo del NNA.
		\item \textbf{Estatus escolar:} Situación escolar actual del NNA.
		\item \textbf{Equipo asignado:} Nombre del equipo multidisciplinario asignado al NNA.
		\item \textbf{Estatus:} Estado actual del registro del NNA en el sistema.
		\item \textbf{Acciones:} Botón \IUbutton{Editar} por cada fila.
		\item \textbf{Gestión:} Botón \IUbutton{Eliminar} por cada fila.
	\end{itemize}
	\item Tabla de Tutores con los siguientes campos por registro:
	\begin{itemize}
		\item \textbf{Nombre:} Nombre completo del tutor.
		\item \textbf{CURP:} Clave Única de Registro de Población del tutor.
		\item \textbf{Parentesco:} Relación del tutor con el NNA.
		\item \textbf{Sexo:} Sexo del tutor.
	\end{itemize}
	\item Mensaje \MSGref{MSG21}{No hay NNAs registrados} si no existen NNAs en el sistema.
	\item Mensaje \MSGref{MSG22}{No hay tutores registrados} si no existen tutores en el sistema.
\end{itemize}

\subsection{Entradas}
\begin{itemize}
	\item No requiere entradas directas del usuario. Ambos listados se cargan automáticamente al acceder a la pantalla.
\end{itemize}

\subsection{Comandos}
\begin{itemize}
	\item \IUbutton{Editar}: Disponible por cada fila de la tabla de NNAs en la columna Acciones. Redirige a la pantalla de modificación del NNA seleccionado.
	\item \IUbutton{Eliminar}: Disponible por cada fila de la tabla de NNAs en la columna Gestión. Elimina el registro del NNA seleccionado del sistema y actualiza la tabla.
	\item \IUbutton{Volver}: Regresa a la \IUref{IU02}{Pantalla del Panel de Administrador}.
\end{itemize}


% !TeX root = ../proyecto.tex
%--------------------------------------
\section{IU12 Pantalla de Registro de NNA}

\subsection{Objetivo}
Permitir a la directora y coordinadora registrar en el sistema SICORRE los datos personales, escolares y de domicilio de una Niña, Niño o Adolescente (NNA), así como asignarle un tutor y un equipo multidisciplinario.

\subsection{Diseño}
Esta pantalla \IUref{IU12}{Pantalla de Registro de NNA} es accesible desde la \IUref{IU11}{Pantalla de Consulta de NNAs y Tutores}. El formulario se organiza en tres secciones: \textbf{Datos personales}, \textbf{Escolaridad y asignación} y \textbf{Domicilio}.

La sección de \textbf{Datos personales} contiene campos para ingresar el nombre, apellido paterno, apellido materno, fecha de nacimiento, sexo, CURP (opcional) y nacionalidad del NNA. Adicionalmente, presenta el checkbox \IUbutton{Es migrante} para indicar si el NNA tiene condición migratoria.

La sección de \textbf{Escolaridad y asignación} contiene un selector de estatus escolar, un selector de tutor y un selector de equipo asignado.

La sección de \textbf{Domicilio} permite ingresar el código postal y utilizar el botón \IUbutton{Buscar} para autocompletar los campos de la dirección. Una vez ejecutada la búsqueda, se habilitan los campos de calle, número exterior, número interior, pueblo o comunidad y tipo de vivienda.

\IUfig[.9]{RegistroNNA}{IU12}{Pantalla de Registro de NNA.}

\subsection{Salidas}
\begin{itemize}
	\item Formulario con los campos del NNA organizados en tres secciones:
	\begin{itemize}
		\item \textbf{Datos personales:} Nombre(s), apellido paterno, apellido materno, fecha de nacimiento, sexo, CURP y nacionalidad.
		\item \textbf{Escolaridad y asignación:} Estatus escolar, tutor asignado y equipo asignado.
		\item \textbf{Domicilio:} Código postal, calle, número exterior, número interior, pueblo o comunidad y tipo de vivienda.
	\end{itemize}
	\item Mensaje \MSGref{MSG23}{NNA registrado exitosamente} al guardar el registro correctamente.
	\item Mensaje \MSGref{MSG24}{Error al registrar NNA} si ocurre un problema al guardar el registro.
\end{itemize}

\subsection{Entradas}
\begin{itemize}
	\item \textbf{Nombre(s):} Campo de texto obligatorio. Nombre(s) del NNA.
	\item \textbf{Apellido paterno:} Campo de texto obligatorio.
	\item \textbf{Apellido materno:} Campo de texto obligatorio.
	\item \textbf{Fecha de nacimiento:} Campo de fecha obligatorio con formato dd/mm/aaaa.
	\item \textbf{Sexo:} Selector desplegable obligatorio.
	\item \textbf{CURP:} Campo de texto opcional.
	\item \textbf{Nacionalidad:} Campo de texto obligatorio. Valor por defecto: \textit{Mexicana}.
	\item \textbf{Es migrante:} Checkbox opcional. Indica si el NNA tiene condición migratoria.
	\item \textbf{Estatus escolar:} Selector desplegable obligatorio.
	\item \textbf{Seleccionar tutor:} Selector desplegable obligatorio.
	\item \textbf{Equipo asignado:} Selector desplegable obligatorio.
	\item \textbf{Código postal:} Campo de texto obligatorio. Se utiliza junto con el botón \IUbutton{Buscar} para autocompletar los campos de domicilio.
	\item \textbf{Calle:} Campo de texto obligatorio. Se autocompleta al buscar por código postal.
	\item \textbf{Núm. exterior:} Campo de texto obligatorio.
	\item \textbf{Núm. interior:} Campo de texto opcional.
	\item \textbf{Pueblo / Comunidad:} Campo de texto obligatorio. Se autocompleta al buscar por código postal.
	\item \textbf{Tipo de vivienda:} Selector desplegable obligatorio.
\end{itemize}

\subsection{Comandos}
\begin{itemize}
	\item \IUbutton{Buscar}: Localiza el código postal ingresado y autocompleta los campos de domicilio correspondientes.
	\item \IUbutton{Guardar registro}: Valida los campos del formulario y almacena el registro del NNA en el sistema. Muestra el mensaje \MSGref{MSG23}{NNA registrado exitosamente} si el proceso concluye correctamente.
	\item \IUbutton{Cancelar}: Descarta los datos ingresados y regresa a la \IUref{IU11}{Pantalla de Consulta de NNAs y Tutores} sin realizar cambios.
\end{itemize}





% !TeX root = ../proyecto.tex
%--------------------------------------
\section{IU13 Pantalla de Registro de Tutor}

\subsection{Objetivo}
Permitir a la directora y coordinadora registrar en el sistema SICORRE los datos personales de un tutor vinculado a una Niña, Niño o Adolescente (NNA).

\subsection{Diseño}
Esta pantalla \IUref{IU13}{Pantalla de Registro de Tutor} es accesible desde el botón \IUbutton{Registrar Tutor} ubicado en la sección de tutores de la \IUref{IU11}{Pantalla de Consulta de NNAs y Tutores}. Presenta un formulario con los datos personales del tutor: nombre, primer apellido, segundo apellido, fecha de nacimiento, sexo, CURP, parentesco y nacionalidad. Adicionalmente, incluye el checkbox \IUbutton{Es migrante} para indicar si el tutor tiene condición migratoria.

{\tiny }\IUfig[.9]{RegistroTutor}{IU13}{Pantalla de Registro de Tutor.}

\subsection{Salidas}
\begin{itemize}
	\item Formulario con los datos personales del tutor.
	\item Mensaje \MSGref{MSG25}{Tutor registrado exitosamente} al guardar el registro correctamente.
	\item Mensaje \MSGref{MSG26}{Error al registrar tutor} si ocurre un problema al guardar el registro.
\end{itemize}

\subsection{Entradas}
\begin{itemize}
	\item \textbf{Nombre:} Campo de texto obligatorio. Nombre(s) del tutor.
	\item \textbf{Primer apellido:} Campo de texto obligatorio.
	\item \textbf{Segundo apellido:} Campo de texto obligatorio.
	\item \textbf{Fecha de nacimiento:} Campo de fecha obligatorio con formato dd/mm/aaaa.
	\item \textbf{Sexo:} Selector desplegable obligatorio.
	\item \textbf{CURP:} Campo de texto opcional.
	\item \textbf{Parentesco:} Campo de texto obligatorio. Indica la relación del tutor con el NNA (Ej. Padre, Madre).
	\item \textbf{Nacionalidad:} Campo de texto obligatorio. Valor por defecto: \textit{Mexicana}.
	\item \textbf{Es migrante:} Checkbox opcional. Indica si el tutor tiene condición migratoria.
\end{itemize}

\subsection{Comandos}
\begin{itemize}
	\item \IUbutton{Guardar registro}: Valida los campos del formulario y almacena el registro del tutor en el sistema. Muestra el mensaje \MSGref{MSG25}{Tutor registrado exitosamente} si el proceso concluye correctamente.
	\item \IUbutton{Cancelar}: Descarta los datos ingresados y regresa a la \IUref{IU11}{Pantalla de Consulta de NNAs y Tutores} sin realizar cambios.
\end{itemize}

% !TeX root = ../proyecto.tex
%--------------------------------------
\section{IU14 Pantalla de Consulta de NNA}

\subsection{Objetivo}
Permitir a la directora y coordinadora consultar el perfil completo de una Niña, 
Niño o Adolescente (NNA) registrado en el sistema SICORRE, incluyendo su 
información personal, tutor, equipo asignado, estatus escolar y domicilio.

\subsection{Diseño}
Esta pantalla \IUref{IU14}{Pantalla de Consulta de NNA} es accesible al 
seleccionar el nombre de un NNA en la tabla correspondiente de la 
\IUref{IU11}{Pantalla de Consulta de NNAs y Tutores}. Muestra el perfil completo 
del NNA organizado en las siguientes secciones: \textbf{Información Personal}, 
\textbf{Tutor}, \textbf{Equipo}, \textbf{Estatus Escolar} y \textbf{Dirección}. 
La pantalla también presenta una fotografía del NNA en el encabezado junto con 
su nombre completo.

\IUfig[.9]{ConsultaNNA}{IU14}{Pantalla de Consulta de NNA.}

\subsection{Salidas}
\begin{itemize}
	\item Fotografía y nombre completo del NNA en el encabezado.
	\item \textbf{Información Personal:}
	\begin{itemize}
		\item \textbf{Fecha de nacimiento:} Fecha de nacimiento del NNA.
		\item \textbf{Sexo:} Sexo del NNA.
		\item \textbf{CURP:} Clave Única de Registro de Población del NNA.
	\end{itemize}
	\item \textbf{Tutor:} Nombre completo del tutor asignado al NNA.
	\item \textbf{Equipo:} Nombre del equipo multidisciplinario asignado, o el 
	texto \textit{No asignado} si el NNA no cuenta con equipo.
	\item \textbf{Estatus Escolar:} Situación escolar actual del NNA.
	\item \textbf{Dirección:}
	\begin{itemize}
		\item \textbf{Calle:} Nombre de la calle del domicilio.
		\item \textbf{No. Exterior:} Número exterior del domicilio.
		\item \textbf{No. Interior:} Número interior del domicilio.
		\item \textbf{Colonia:} Colonia del domicilio.
		\item \textbf{Código Postal:} Código postal del domicilio.
		\item \textbf{Pueblo/Comunidad:} Pueblo o comunidad del domicilio, o el 
		texto \textit{No especificado} si no fue proporcionado.
		\item \textbf{Tipo de vivienda:} Tipo de vivienda del domicilio.
	\end{itemize}
\end{itemize}

\subsection{Entradas}
\begin{itemize}
	\item No requiere entradas directas del usuario. La información se carga 
	automáticamente al seleccionar el NNA desde la 
	\IUref{IU11}{Pantalla de Consulta de NNAs y Tutores}.
\end{itemize}

\subsection{Comandos}
\begin{itemize}
	\item \IUbutton{Modificar}: Redirige a la 
	\IUref{IU15}{Pantalla de Modificación de NNA} para editar los datos 
	del NNA actual.
	\item \IUbutton{Eliminar}: Elimina el registro del NNA del sistema y 
	regresa a la \IUref{IU11}{Pantalla de Consulta de NNAs y Tutores}.
	\item \IUbutton{Volver}: Regresa a la 
	\IUref{IU11}{Pantalla de Consulta de NNAs y Tutores} sin realizar cambios.
\end{itemize}


% !TeX root = ../proyecto.tex
%--------------------------------------
\section{IU15 Pantalla de Modificación de NNA}

\subsection{Objetivo}
Permitir a la directora y coordinadora actualizar los datos personales, escolares 
y de domicilio de una Niña, Niño o Adolescente (NNA) previamente registrado en 
el sistema SICORRE.

\subsection{Diseño}
Esta pantalla \IUref{IU15}{Pantalla de Modificación de NNA} es accesible desde 
el botón \IUbutton{Modificar} de la \IUref{IU14}{Pantalla de Consulta de NNA}. 
Presenta el mismo formulario que la \IUref{IU12}{Pantalla de Registro de NNA}, 
organizado en tres secciones: \textbf{Datos personales}, \textbf{Escolaridad y 
	asignación} y \textbf{Domicilio}. Todos los campos se cargan automáticamente con 
la información actual del NNA, permitiendo al usuario modificar únicamente los 
datos que requieran actualización.

\IUfig[.9]{ModificacionNNA}{IU15}{Pantalla de Modificación de NNA.}

\subsection{Salidas}
\begin{itemize}
	\item Formulario con los campos del NNA precargados con su información actual, 
	organizados en tres secciones:
	\begin{itemize}
		\item \textbf{Datos personales:} Nombre(s), apellido paterno, apellido 
		materno, fecha de nacimiento, sexo, CURP y nacionalidad.
		\item \textbf{Escolaridad y asignación:} Estatus escolar, tutor asignado 
		y equipo asignado.
		\item \textbf{Domicilio:} Código postal, calle, número exterior, número 
		interior, pueblo o comunidad y tipo de vivienda.
	\end{itemize}
	\item Mensaje \MSGref{MSG25}{NNA actualizado exitosamente} al guardar los 
	cambios correctamente.
	\item Mensaje \MSGref{MSG26}{Error al actualizar NNA} si ocurre un problema 
	al guardar los cambios.
\end{itemize}

\subsection{Entradas}
\begin{itemize}
	\item \textbf{Nombre(s):} Campo de texto obligatorio. Precargado con el 
	nombre actual del NNA.
	\item \textbf{Apellido paterno:} Campo de texto obligatorio. Precargado con 
	el valor actual.
	\item \textbf{Apellido materno:} Campo de texto obligatorio. Precargado con 
	el valor actual.
	\item \textbf{Fecha de nacimiento:} Campo de fecha obligatorio con formato 
	dd/mm/aaaa. Precargado con el valor actual.
	\item \textbf{Sexo:} Selector desplegable obligatorio. Precargado con el 
	valor actual.
	\item \textbf{CURP:} Campo de texto opcional. Precargado con el valor actual.
	\item \textbf{Nacionalidad:} Campo de texto obligatorio. Precargado con el 
	valor actual.
	\item \textbf{Es migrante:} Checkbox opcional. Precargado con el estado 
	actual del NNA.
	\item \textbf{Estatus escolar:} Selector desplegable obligatorio. Precargado 
	con el valor actual.
	\item \textbf{Seleccionar tutor:} Selector desplegable obligatorio. Precargado 
	con el tutor actualmente asignado.
	\item \textbf{Equipo asignado:} Selector desplegable obligatorio. Precargado 
	con el equipo actualmente asignado.
	\item \textbf{Código postal:} Campo de texto obligatorio. Precargado con el 
	valor actual.
	\item \textbf{Calle:} Campo de texto obligatorio. Precargado con el valor 
	actual.
	\item \textbf{Núm. exterior:} Campo de texto obligatorio. Precargado con el 
	valor actual.
	\item \textbf{Núm. interior:} Campo de texto opcional. Precargado con el 
	valor actual.
	\item \textbf{Pueblo / Comunidad:} Campo de texto obligatorio. Precargado 
	con el valor actual.
	\item \textbf{Tipo de vivienda:} Selector desplegable obligatorio. Precargado 
	con el valor actual.
\end{itemize}

\subsection{Comandos}
\begin{itemize}
	\item \IUbutton{Buscar}: Localiza el código postal ingresado y autocompleta 
	los campos de domicilio correspondientes.
	\item \IUbutton{Actualizar}: Valida los campos del formulario y almacena los 
	cambios del NNA en el sistema. Muestra el mensaje 
	\MSGref{MSG25}{NNA actualizado exitosamente} si el proceso concluye 
	correctamente.
	\item \IUbutton{Cancelar}: Descarta los cambios realizados y regresa a la 
	\IUref{IU14}{Pantalla de Consulta de NNA} sin modificar ningún dato.
\end{itemize}


% !TeX root = ../proyecto.tex
%--------------------------------------
\section{IU16 Pantalla del Panel de Restitución de Derechos}

\subsection{Objetivo}
Permitir a la directora y coordinadora acceder a las funcionalidades relacionadas 
con la restitución de derechos en el sistema SICORRE, mediante la navegación hacia 
la consulta de catálogos y la consulta de actores en materia de derechos.

\subsection{Diseño}
Esta pantalla \IUref{IU16}{Pantalla del Panel de Restitución de Derechos} es 
accesible desde la \IUref{IU02}{Pantalla del Panel de Administrador}. Presenta 
dos botones de navegación: \IUbutton{Consultar Catálogos} y 
\IUbutton{Consultar Actores en Materia de Derechos}. Actualmente únicamente se 
encuentra habilitado el botón \IUbutton{Consultar Catálogos}.

\IUfig[.9]{PanelRestitucion}{IU16}{Pantalla del Panel de Restitución de Derechos.}

\subsection{Salidas}
\begin{itemize}
	\item Botón \IUbutton{Consultar Catálogos} habilitado para acceder a la 
	consulta de discapacidades, enfermedades y domicilios.
	\item Botón \IUbutton{Consultar Actores en Materia de Derechos} deshabilitado 
	en la iteración actual del sistema.
\end{itemize}

\subsection{Entradas}
\begin{itemize}
	\item No requiere entradas del usuario.
\end{itemize}

\subsection{Comandos}
\begin{itemize}
	\item \IUbutton{Consultar Catálogos}: Redirige a la 
	\IUref{IU17}{Pantalla de Consulta de Catálogos}.
	\item \IUbutton{Consultar Actores en Materia de Derechos}: No disponible en 
	la iteración actual del sistema.
	\item \IUbutton{Volver}: Regresa a la \IUref{IU02}{Pantalla del Panel de 
		Administrador}.
\end{itemize}

% !TeX root = ../proyecto.tex
%--------------------------------------
\section{IU17 Pantalla de Consulta de Catálogos}

\subsection{Objetivo}
Permitir a la directora y coordinadora consultar los catálogos utilizados en los 
formularios del sistema SICORRE, mostrando los valores disponibles para cada 
categoría.

\subsection{Diseño}
Esta pantalla \IUref{IU17}{Pantalla de Consulta de Catálogos} es accesible desde 
el botón \IUbutton{Consultar Catálogos} de la \IUref{IU16}{Pantalla del Panel de 
	Restitución de Derechos}. Presenta una serie de tablas, una por cada catálogo del 
sistema, cada una encabezada por el nombre del catálogo y el número de registros 
disponibles. Cada tabla muestra dos columnas: \textbf{ID} y \textbf{Descripción}. 
En caso de que un catálogo no cuente con registros, se muestra el mensaje 
\MSGref{MSG27}{Sin registros disponibles} en lugar de la tabla.

Los catálogos presentados son: \textbf{Tipos de vivienda}, \textbf{Vivienda NNA}, 
\textbf{Estatus escolar}, \textbf{Idiomas}, \textbf{Enfermedades} y 
\textbf{Discapacidades}.

\IUfig[.9]{ConsultaCatalogos}{IU17}{Pantalla de Consulta de Catálogos.}

\subsection{Salidas}
\begin{itemize}
	\item Serie de tablas por cada catálogo del sistema, mostrando para cada una:
	\begin{itemize}
		\item \textbf{Nombre del catálogo} y número de registros en el encabezado.
		\item \textbf{ID:} Identificador numérico del registro.
		\item \textbf{Descripción:} Valor descriptivo del registro.
	\end{itemize}
	\item Mensaje \MSGref{MSG27}{Sin registros disponibles} en los catálogos que 
	no cuenten con registros.
\end{itemize}

\subsection{Entradas}
\begin{itemize}
	\item No requiere entradas del usuario. Los catálogos se cargan 
	automáticamente al acceder a la pantalla.
\end{itemize}

\subsection{Comandos}
\begin{itemize}
	\item \IUbutton{Volver al panel}: Regresa a la 
	\IUref{IU16}{Pantalla del Panel de Restitución de Derechos}.
\end{itemize}


%!TEX root = ejemplo.tex
%=========================================================
\section{Catálogo de mensajes}	
\label{sec:mensajes}

En esta sección se describen todos los mensajes que aparecen en el sistema. Para cada mensaje se especifica:

\begin{description}\itemsep0em
	\item[Id:] Identificador del mensaje de la forma ``MSG XX'' y descripción corta del mismo.
	\item[Tipo:] Tipo del mensaje el cual puede ser: 
	\begin{description}
		\item[Normal:] Mensaje que informa al usuario una instrucción o el estado interno que guarda el sistema, suele tener un color {\color{msgNormalColor}Azul}.
		\item[Éxito:] Mensaje que informa al usuario sobre una acción realizada, sirve para confirmar el correcto funcionamiento del sistema. Se presentan con un color {\color{msgInfoColor}Verde}.
		\item[Atención:] Mensaje que tiene como finalidad llamar la atención del usuario a una situación que requiere su intervención, por ejemplo cuando una actividad ha generado un efecto colateral o se realizará una acción destructiva y no reversible. Se presentan con un color {\color{msgWarningColor}Naranja}.
		\item[Error:] Mensaje que informa al usuario un fallo en en una operación o un impedimento para realizarla, por ejemplo: cuando no se puede efectuar la acción solicitada, cuando un dato falta o tiene un formato no aceptado por el sistema. Se presentan con un color {\color{msgErrorColor}Rojo}.
	\end{description}
	\item[Propósito:] Explicación del propósito del mensaje.
	\item[Redacción:] Redacción del mensaje.
	\item[Parámetros:] En caso de que el mensaje pueda variar se especifican los casos y la forma en que debe adaptarse la redacción
	\item[Ejemplos:] Ejemplos de como debe renderizarse el mensaje.
\end{description}


\subsection{Lista de mensajes}

%--------------------------------------
% MENSAJES DE CONTROL DE ACCESO
%--------------------------------------

\begin{cdtMessage}[msgInfoColor]{MSG01}{Bienvenida al usuario}
	\item[Propósito:] Indicar al usuario que ha ingresado 
	satisfactoriamente al sistema SICORRE.
	\item[Redacción:] Bienvenido/a $<$nombre$>$.
	\item[Parámetros:] \hspace{1cm}
	\begin{itemize}
		\item $<$nombre$>$ Nombre completo del usuario autenticado.
	\end{itemize}
	\item[Ejemplos:] Bienvenido/a Joji Pérez Ramírez.
\end{cdtMessage}

\begin{cdtMessage}[msgErrorColor]{MSG02}{Credenciales incorrectas}
	\item[Propósito:] Indicar al usuario que el correo electrónico 
	o la contraseña ingresados no son válidos.
	\item[Redacción:] El correo electrónico o la contraseña 
	son incorrectos. Verifique sus datos e intente de nuevo.
	\item[Parámetros:] \hspace{1cm}
	\begin{itemize}
		\item Ninguno. Por seguridad no se especifica cuál de 
		los dos campos es incorrecto.
	\end{itemize}
	\item[Ejemplos:] El correo electrónico o la contraseña 
	son incorrectos. Verifique sus datos e intente de nuevo.
\end{cdtMessage}

\begin{cdtMessage}[msgErrorColor]{MSG03}{Acceso revocado}
	\item[Propósito:] Indicar al usuario que su cuenta ha sido 
	revocada y no puede acceder al sistema.
	\item[Redacción:] El acceso de $<$nombre$>$ ha sido revocado. 
	Contacte al administrador para más información.
	\item[Parámetros:] \hspace{1cm}
	\begin{itemize}
		\item $<$nombre$>$ Nombre completo del usuario 
		que intenta acceder.
	\end{itemize}
	\item[Ejemplos:] El acceso de Joji Pérez Ramírez ha sido 
	revocado. Contacte al administrador para más información.
\end{cdtMessage}

%--------------------------------------
% MENSAJES DE GESTIÓN DE USUARIOS
%--------------------------------------

\begin{cdtMessage}[msgWarningColor]{MSG04}{Confirmar revocación de acceso}
	\item[Propósito:] Solicitar confirmación al administrador 
	antes de revocar el acceso de un usuario al sistema.
	\item[Redacción:] ¿Está seguro que desea revocar el acceso 
	de $<$nombre$>$? El usuario no podrá ingresar al sistema, 
	pero su historial de actividad se conservará.
	\item[Parámetros:] \hspace{1cm}
	\begin{itemize}
		\item $<$nombre$>$ Nombre completo del usuario 
		a quien se revocará el acceso.
	\end{itemize}
	\item[Ejemplos:] ¿Está seguro que desea revocar el acceso 
	de Joji Pérez Ramírez? El usuario no podrá ingresar al 
	sistema, pero su historial de actividad se conservará.
\end{cdtMessage}

\begin{cdtMessage}[msgWarningColor]{MSG05}{Confirmar eliminación de usuario}
	\item[Propósito:] Solicitar confirmación al administrador 
	antes de eliminar permanentemente el registro de un usuario.
	\item[Redacción:] ¿Está seguro que desea eliminar el registro 
	de $<$nombre$>$? Esta acción no puede deshacerse.
	\item[Parámetros:] \hspace{1cm}
	\begin{itemize}
		\item $<$nombre$>$ Nombre completo del usuario 
		a eliminar del sistema.
	\end{itemize}
	\item[Ejemplos:] ¿Está seguro que desea eliminar el registro 
	de Joji Pérez Ramírez? Esta acción no puede deshacerse.
\end{cdtMessage}

\begin{cdtMessage}[msgInfoColor]{MSG06}{Registro exitoso de usuario}
	\item[Propósito:] Confirmar al administrador que el nuevo 
	usuario ha sido registrado correctamente en el sistema.
	\item[Redacción:] El usuario $<$nombre$>$ ha sido registrado 
	exitosamente en el sistema.
	\item[Parámetros:] \hspace{1cm}
	\begin{itemize}
		\item $<$nombre$>$ Nombre completo del usuario 
		recién registrado.
	\end{itemize}
	\item[Ejemplos:] El usuario Joji Pérez Ramírez ha sido 
	registrado exitosamente en el sistema.
\end{cdtMessage}

\begin{cdtMessage}[msgErrorColor]{MSG07}{CURP ya registrada}
	\item[Propósito:] Indicar al administrador que la CURP 
	ingresada ya está asociada a otro usuario en el sistema.
	\item[Redacción:] La CURP $<$curp$>$ ya se encuentra 
	registrada en el sistema.
	\item[Parámetros:] \hspace{1cm}
	\begin{itemize}
		\item $<$curp$>$ CURP ingresada en el formulario.
	\end{itemize}
	\item[Ejemplos:] La CURP PERJ900101HDFRZN09 ya se encuentra 
	registrada en el sistema.
\end{cdtMessage}

\begin{cdtMessage}[msgErrorColor]{MSG08}{Correo ya registrado}
	\item[Propósito:] Indicar al administrador que el correo 
	electrónico ingresado ya está asociado a otro usuario 
	en el sistema.
	\item[Redacción:] El correo $<$correo$>$ ya se encuentra 
	registrado en el sistema.
	\item[Parámetros:] \hspace{1cm}
	\begin{itemize}
		\item $<$correo$>$ Correo electrónico ingresado 
		en el formulario.
	\end{itemize}
	\item[Ejemplos:] El correo joji.perez@fundacion.org ya 
	se encuentra registrado en el sistema.
\end{cdtMessage}

\begin{cdtMessage}[msgErrorColor]{MSG09}{Código postal no encontrado}
	\item[Propósito:] Indicar al usuario que el código postal 
	ingresado no existe en el catálogo SEPOMEX del sistema.
	\item[Redacción:] El código postal $<$cp$>$ no fue encontrado. 
	Verifique e intente de nuevo.
	\item[Parámetros:] \hspace{1cm}
	\begin{itemize}
		\item $<$cp$>$ Código postal ingresado en el formulario.
	\end{itemize}
	\item[Ejemplos:] El código postal 00000 no fue encontrado. 
	Verifique e intente de nuevo.
\end{cdtMessage}

\begin{cdtMessage}[msgErrorColor]{MSG10}{Campos obligatorios incompletos}
	\item[Propósito:] Indicar al usuario que existen campos 
	obligatorios sin completar antes de guardar el formulario.
	\item[Redacción:] El campo $<$campo$>$ es obligatorio 
	y no puede estar vacío.
	\item[Parámetros:] \hspace{1cm}
	\begin{itemize}
		\item $<$campo$>$ Nombre del campo obligatorio 
		que se encuentra vacío.
	\end{itemize}
	\item[Ejemplos:] El campo Nombre(s) es obligatorio 
	y no puede estar vacío.
\end{cdtMessage}

\begin{cdtMessage}[msgInfoColor]{MSG11}{Modificación exitosa de usuario}
	\item[Propósito:] Confirmar al administrador que los datos 
	del usuario han sido actualizados correctamente en el sistema.
	\item[Redacción:] Los datos de $<$nombre$>$ han sido 
	actualizados exitosamente.
	\item[Parámetros:] \hspace{1cm}
	\begin{itemize}
		\item $<$nombre$>$ Nombre completo del usuario modificado.
	\end{itemize}
	\item[Ejemplos:] Los datos de Joji Pérez Ramírez han sido 
	actualizados exitosamente.
\end{cdtMessage}

%--------------------------------------
% MENSAJES DE GESTIÓN DE EQUIPOS
%--------------------------------------

\begin{cdtMessage}[msgWarningColor]{MSG12}{Sin equipos registrados}
	\item[Propósito:] Informar al administrador que no existe 
	ningún equipo multidisciplinario registrado en el sistema.
	\item[Redacción:] No hay equipos multidisciplinarios 
	registrados en el sistema.
	\item[Parámetros:] \hspace{1cm}
	\begin{itemize}
		\item Ninguno.
	\end{itemize}
	\item[Ejemplos:] No hay equipos multidisciplinarios 
	registrados en el sistema.
\end{cdtMessage}

\begin{cdtMessage}[msgWarningColor]{MSG13}{Confirmar eliminación de equipo}
	\item[Propósito:] Solicitar confirmación al administrador 
	antes de eliminar permanentemente el registro de un equipo 
	multidisciplinario.
	\item[Redacción:] ¿Está seguro que desea eliminar el equipo 
	$<$nombre$>$? Esta acción no puede deshacerse.
	\item[Parámetros:] \hspace{1cm}
	\begin{itemize}
		\item $<$nombre$>$ Nombre del equipo a eliminar.
	\end{itemize}
	\item[Ejemplos:] ¿Está seguro que desea eliminar el equipo 
	Equipo Alpha? Esta acción no puede deshacerse.
\end{cdtMessage}

\begin{cdtMessage}[msgInfoColor]{MSG14}{Equipo registrado exitosamente}
	\item[Propósito:] Confirmar al administrador que el nuevo 
	equipo multidisciplinario ha sido registrado correctamente 
	en el sistema.
	\item[Redacción:] El equipo $<$nombre$>$ ha sido registrado 
	exitosamente en el sistema.
	\item[Parámetros:] \hspace{1cm}
	\begin{itemize}
		\item $<$nombre$>$ Nombre del equipo recién registrado.
	\end{itemize}
	\item[Ejemplos:] El equipo Equipo Alpha ha sido registrado 
	exitosamente en el sistema.
\end{cdtMessage}

\begin{cdtMessage}[msgErrorColor]{MSG15}{Nombre de equipo ya registrado}
	\item[Propósito:] Indicar al administrador que el nombre 
	ingresado para el equipo ya existe en el sistema.
	\item[Redacción:] El nombre de equipo $<$nombre$>$ ya se 
	encuentra registrado en el sistema.
	\item[Parámetros:] \hspace{1cm}
	\begin{itemize}
		\item $<$nombre$>$ Nombre de equipo ingresado 
		en el formulario.
	\end{itemize}
	\item[Ejemplos:] El nombre de equipo Equipo Alpha ya se 
	encuentra registrado en el sistema.
\end{cdtMessage}

\begin{cdtMessage}[msgWarningColor]{MSG16}{No hay integrantes disponibles}
	\item[Propósito:] Informar al administrador que todo el 
	personal activo ya pertenece a un equipo multidisciplinario 
	y no hay integrantes disponibles para asignar.
	\item[Redacción:] No hay personal disponible para agregar 
	al equipo. Todos los integrantes activos ya pertenecen 
	a un equipo multidisciplinario.
	\item[Parámetros:] \hspace{1cm}
	\begin{itemize}
		\item Ninguno.
	\end{itemize}
	\item[Ejemplos:] No hay personal disponible para agregar 
	al equipo. Todos los integrantes activos ya pertenecen 
	a un equipo multidisciplinario.
\end{cdtMessage}

\begin{cdtMessage}[msgInfoColor]{MSG17}{Integrante agregado exitosamente}
	\item[Propósito:] Confirmar al administrador que el integrante 
	seleccionado ha sido asignado correctamente al equipo 
	multidisciplinario.
	\item[Redacción:] $<$nombre$>$ ha sido agregado exitosamente 
	al equipo $<$equipo$>$.
	\item[Parámetros:] \hspace{1cm}
	\begin{itemize}
		\item $<$nombre$>$ Nombre completo del integrante agregado.
		\item $<$equipo$>$ Nombre del equipo al que fue asignado.
	\end{itemize}
	\item[Ejemplos:] Joji Pérez Ramírez ha sido agregado 
	exitosamente al equipo Equipo Alpha.
\end{cdtMessage}

%--------------------------------------
% MENSAJES DE CONSULTA GENERAL
%--------------------------------------

\begin{cdtMessage}[msgWarningColor]{MSG18}{Sin usuarios registrados}
	\item[Propósito:] Informar al administrador que no existe 
	ningún usuario registrado en el sistema.
	\item[Redacción:] No hay usuarios registrados en el sistema.
	\item[Parámetros:] \hspace{1cm}
	\begin{itemize}
		\item Ninguno.
	\end{itemize}
	\item[Ejemplos:] No hay usuarios registrados en el sistema.
\end{cdtMessage}
\begin{cdtMessage}[msgInfoColor]{MSG19}{Equipo modificado exitosamente}
	\begin{description}\itemsep0em
		\item[Descripción:] El nombre del equipo fue actualizado correctamente en el sistema.
	\end{description}
\end{cdtMessage}

\begin{cdtMessage}[msgInfoColor]{MSG20}{Integrante eliminado exitosamente}
	\begin{description}\itemsep0em
		\item[Descripción:] El integrante fue removido del equipo correctamente.
	\end{description}
\end{cdtMessage}

\begin{cdtMessage}[msgWarningColor]{MSG21}{No hay NNAs registrados}
	\begin{description}\itemsep0em
		\item[Descripción:] No existen NNAs registrados en el sistema.
	\end{description}
\end{cdtMessage}

\begin{cdtMessage}[msgWarningColor]{MSG22}{No hay tutores registrados}
	\begin{description}\itemsep0em
		\item[Descripción:] No existen tutores registrados en el sistema.
	\end{description}
\end{cdtMessage}

\begin{cdtMessage}[msgInfoColor]{MSG23}{NNA registrado exitosamente}
	\begin{description}\itemsep0em
		\item[Descripción:] El registro del NNA fue almacenado correctamente en el sistema.
	\end{description}
\end{cdtMessage}

\begin{cdtMessage}[msgErrorColor]{MSG24}{Error al registrar NNA}
	\begin{description}\itemsep0em
		\item[Descripción:] Ocurrió un problema al intentar guardar el registro del NNA.
	\end{description}
\end{cdtMessage}

\begin{cdtMessage}[msgInfoColor]{MSG25}{Tutor registrado exitosamente}
	\begin{description}\itemsep0em
		\item[Descripción:] El registro del tutor fue almacenado correctamente en el sistema.
	\end{description}
\end{cdtMessage}

\begin{cdtMessage}[msgErrorColor]{MSG26}{Error al registrar tutor}
	\begin{description}\itemsep0em
		\item[Descripción:] Ocurrió un problema al intentar guardar el registro del tutor.
	\end{description}
\end{cdtMessage}

\begin{cdtMessage}[msgWarningColor]{MSG27}{Sin registros disponibles}
	\begin{description}\itemsep0em
		\item[Descripción:] No hay registros disponibles para mostrar en el catálogo.
	\end{description}
\end{cdtMessage}
